% Source: https://artofproblemsolving.com/community/c3h2355672_an_interesting_problem
% Stable identity: aops-topic-2355672-q01
\textbf{(a)} Fix $k$. For each starting position $j=1,2,3,\dots$ look at the block of $k+1$ digits
\[\left(x_j,x_{j+1},\dots,x_{j+k}\right).\]
There are only $10^{k+1}$ possible blocks of $k+1$ digits, but infinitely many starting positions. By the pigeonhole principle, among the starting positions $1,2,\dots,10^{k+1}+1$ two different ones, $m\ne n$, give the same block, which is exactly the required condition. Since $k$ was arbitrary, $x$ is repetitive. $\blacksquare$

\textbf{(b)} If $x$ is repetitive, then $0\le x<1$, so $0\le x^2<1$ as well, and by (a) $x^2$ is repetitive. $\blacksquare$

(With the definition as stated, the property holds for every number, so the claim about squares is immediate; it would only be a real restriction if, for example, the two blocks were required to be adjacent.)
